%%%%%%%%%%%%%%%%%%%%%%%%%%%%%%%%%
%
%       EXERCÍCIO
%
%%%%%%%%%%%%%%%%%%%%%%%%%%%%%%%%%

\ifdefstring{\atividade}{prova}{%
    \renewcommand{\valorquestao}{\ValorQ\ pontos}
}%

\begin{exercicioBanco}[\valorquestao]
    A função de probabilidade da variável \(X\) é \(P(X=k) = \dfrac{1}{5}\) para \(k = 1, 3, \dots, 5\). Calcule \(E(X)\) e \(E(X^{2})\) e, usando esses resultados, determine \(E\big((X+3)^{2}\big)\) e \(V(3X-2)\).
\end{exercicioBanco}

